\section{Irreducibility and automorphy on the seed field}

% BEGIN G65I1
\paragraph{G65I1. Irreducibility on the full cyclotomic kernel ($\mathrm{IrrOf4}$).}
Fix a source datum $d:\mathrm{Dat}$, an acceptable $T\ge T_0$, and
write $E=\overline{\mathbb Q}_\ell$, $F=T.F$ and $L=T.L=F\,T.K$.
Inside $G_{\mathbb Q}=\operatorname{Gal}(\Omega/\mathbb Q)$ use
$G_M=\operatorname{Gal}(\Omega/M)$ and the full cyclotomic character
$\varepsilon$. The tower subgroup is $G'=G_L\cap\ker\varepsilon$.
Set $U=G_{L K'}$ and $H=U\cap\ker\varepsilon=G'\cap G_{K'}$.
The datum supplies a continuous rank-two representation $\rho$,
irreducible over $E$ on each open subgroup, and an integral model
$\rho_1$ isomorphic to it over $E$.

We first show that the restriction to $H$ is irreducible. The group
$U$ is open, $H$ is normal in $U$, and $U/H$ is abelian. If there
are exactly one or two $H$-invariant lines, $U$ permutes them.
The stabilizer of one such line is closed by continuity (stabilizing
a given line is a closed matrix condition) and has finite index in
$U$. It is open, contradicting the hypothesis on $\rho$. If there
are three or more invariant lines, every $\rho(h)$ is scalar,
since a nonscalar two-dimensional matrix has at most two eigenlines.
The projective image of $U$ is therefore abelian. If $\rho(U)$
contains a nonscalar matrix $A$, its one or two eigenlines over $E$
are permuted by $U$, because each $\rho(u)A\rho(u)^{-1}$ is a
scalar multiple of $A$. The same open stabilizer again contradicts
irreducibility. If all matrices on $U$ are scalar, every line is
$U$-invariant. These cases exhaust all ways $H$ could preserve a line.

The second integral factor $\tau$ preserves $E v_+$ and $E v_-$ on
$G_{K'}$, where $v_\pm=(1,\pm1)^t$. Its on-coset matrix has rows
$(a,b),(b,a)$ and acts by $\kappa_\pm=a\pm b$. Its defining
orthogonality gives $\kappa_+\kappa_-=a^2-b^2=1$.
Off the subgroup the rows are $(a,b),(-b,-a)$, so the two lines
are exchanged. E65T3 supplies $t\in G'\setminus G_{K'}$.
The distinguished inertia element $s_0$ at the chosen place over
$\ell$ has $\varepsilon(s_0)$ of infinite order. Choose $n>0$ with
$h=s_0^n\in U$ by the finite index of $U$. Since $h$ still lies in
this inertia group, the prescribed local sign in $\mathrm{TauDat}$
gives $\kappa_+(h)=\varepsilon(h)^m$ or
$\kappa_-(h)=\varepsilon(h)^m$, with $m\ge1$. Thus
$\lambda=\kappa_+(h)$ has infinite order. The commutator
$u=hth^{-1}t^{-1}$ fixes $L$ and lies in both $\ker\varepsilon$
and $G_{K'}$, since the two quotient characters are abelian.
The interchange of lines by $t$ gives
\[
             \kappa_+(u)=\lambda^2,\qquad
             \kappa_-(u)=\lambda^{-2}.
\]
In particular $\kappa_+(u)^2\ne\kappa_-(u)^2$.

Over $H$ the tensor product is the sum of two irreducible
two-dimensional representations $V_\pm=\rho_1|_H\otimes\kappa_\pm$.
Their determinants at $u$ are $\det\rho_1(u)\,\kappa_\pm(u)^2$,
so these representations are nonisomorphic. Any invariant subspace of
their direct sum is a sum of some of the two summands. Indeed, its
intersection with each summand is zero or the whole summand. If both
intersections are zero and the subspace is nonzero, the two projections
would identify it equivariantly with both summands. If it contains one
summand, its projection to the other gives the two remaining possibilities.
Since $t$ exchanges $V_+$ and $V_-$, the representation is irreducible
on $G'$. In the ordered basis
$(e_0\otimes e_0,e_0\otimes e_1,-e_1\otimes e_1,e_1\otimes e_0)$
the tensor matrix is exactly $\mathrm{tn}(\rho_1,\tau)$. The scalar
$\delta(g)$ is invertible, so multiplication by it preserves the same
invariant subspaces. Burnside's matrix-algebra theorem over the
algebraically closed field $E$ now gives the required equality
$\operatorname{span}_E Q_T(G')=M_4(E)$, since this linear span is a
unital algebra under matrix multiplication.
% END G65I1

% BEGIN G65I2
\paragraph{G65I2. The local hypotheses for the automorphy input ($\mathrm{GoodOf4}$).}
This assertion holds already for any good raw datum $D$ and any
$T\ge T_0$. Assume that all rational decomposition
groups over $p\in S$ fix $K'$, and the residual root $T_0.\mathrm{rb}(w_\alpha)$
takes value $-1$ on elements of $G_{F_0}$ outside $G_{K'}$,
where $F_0=T_0.F$. Take a prime $w$ split in the specified sense
$T_0.\mathrm{Splits}(w)$, a character $\theta$ unramified away from
$w$, and a model $\mathrm{Mod4}$ of $\rho$ with the prescribed
$k,g_0,O\subset O_1,\rho_1,x$. Write
$\alpha=\varepsilon\omega^{-1}\xi$. The ratio of the first two
labelled characters is $\alpha$ for either labelling, so
$T_0.\mathrm{rb}(w_\alpha)=\bar\alpha$ on $G_{F_0}$.
The good datum gives $\bar\omega=\bar\varepsilon$ and odd $\ell$.
Consequently $\bar\alpha=\bar\xi$, and any element of $G_{F_0}$
with $\bar\alpha=1$ fixes $K'$.

For $p\in S$, every rational decomposition group above $p$ fixes $K'$
by hypothesis. For $p=w$, the splitting condition says that its
decomposition groups fix $F_0$ and have $\bar r_D=1$. The diagonal
entry indexed by $2$ in the original residual representation is
$\bar\alpha$, so these groups too fix $K'$. At every remaining prime
$p$, we have $p\ne\ell$ because $\ell\in S$. Take any valuation over
$p$ and an element $z$ in its inertia over $L$. Its rational image is
in rational inertia at that valuation and fixes $F_0\subset L$. By
the respective unramifiedness hypotheses,
$\xi(z)=\psi(z)=\theta(z)=1$. Thus $z$ fixes $K'$ and
$\phi(z)=\xi(z)\psi(z)^{\pm1}=1$. The away-from-$\ell$ clause of
$\mathrm{Mod4}$ gives $\rho_1(z)=1$, hence $\rho(z)=1$ by the supplied
isomorphism over $E$. These are precisely the two alternatives in
$\mathrm{GoodAt}(K',L,\rho)$ for each rational prime and every
valuation above it.
% END G65I2

% BEGIN G65I3
\paragraph{G65I3. Automorphy after cancelling the seed twist ($\mathrm{AutOf4}$).}
Return to $d:\mathrm{Dat}$ and acceptable $T\ge T_0$. Write
$R=\mathrm{tenB}(\rho_1,\tau)$ for the untwisted tensor over $E$.
Take the matrix-span hypothesis on $Q_T|_{G'}$ of $\mathrm{AutOf4}$,
which G65I1 supplies. Any $R(G_L)$-invariant subspace is invariant
under $Q_T(G')$ because $Q_T(g)=\delta(g)R(g)$ there. It is then
invariant under all matrices and is zero or the whole space. Thus
$R|_{G_L}$ is irreducible. The element of E65T3 fixes $L$ and moves
$K'$, so $K'\not\subseteq L$. The tower gives a totally complex CM
field $L=F\,T.K$. The field $F$ is abelian Galois over $\mathbb Q$,
as is the quadratic field $T.K$, so their compositum is abelian
Galois by restriction into the product of the two Galois groups.

Apply the supplied $d.\mathrm{aut}=\mathrm{AutOffS}(\iota,K',\rho,R)$
to this field. The local hypotheses are exactly G65I2 applied to
$d.\mathrm{quad}$, $d.w,d.\theta$ and $d.\mathrm{mod}$.
The result is a regular algebraic cuspidal $\pi$ of rank four over
$L$ with $\mathrm{AttC}(\iota,R|_{G_L},\pi)$.
At each $g\in G_L$ the seed representation over $E$ is
$\delta(g)R(g)$ in the very same tensor basis. Twisting on $G_L$
by $\delta^{-1}$ gives $R(g)$ entry by entry. This proves attachment
to the precise representation in $\mathrm{AutOf4}$. The argument
uses the arithmetic automorphy input $d.\mathrm{aut}$ in addition
to the integral matrices and local conditions of $d$.
% END G65I3
